\begin{frame}[t]{\ev{\tagP\,\tagO\,\tagA}Ancestry is a tree and evidence flow is a graph}
\begin{columns}[T,onlytextwidth]
\column{.5\textwidth}
\begin{center}
\begin{tikzpicture}[x=1cm,y=.74cm]
\node[state,minimum width=1.2cm] (s) at (0,2.4) {$S_0$};
\node[state,minimum width=1.2cm] (p1) at (-1.9,1.2) {parent $p_1$};
\node[state,minimum width=1.2cm] (p2) at (1.9,1.2) {parent $p_2$};
\node[candidate,minimum width=1.3cm] (a) at (-1.9,-.3) {child A};
\node[candidate,minimum width=1.3cm] (b) at (1.9,-.3) {child B};
\node[evidence,minimum width=1.6cm] (f) at (0,.75) {feedback\\[-2pt]archive};
\node[evidence,minimum width=1.6cm] (e) at (0,-1.55) {evaluator};
\draw[flow] (s)--(p1);\draw[flow] (s)--(p2);\draw[flow] (p1)--(a);\draw[flow] (p2)--(b);
\draw[message] (f)--(a);\draw[message] (f)--(b);
\draw[-{Latex[length=2mm]},draw=Gold,dashed,line width=.9pt] (a)--(e);
\draw[-{Latex[length=2mm]},draw=Gold,dashed,line width=.9pt] (b)--(e);
\draw[relation] (a.east)--(b.west);
\end{tikzpicture}
\end{center}
\column{.46\textwidth}
{\small
\term{Edge labels}\\[2pt]
\tikz[baseline=-.5ex]\draw[flow](0,0)--(.6,0);\ parent relation: state ancestry\\[2pt]
\tikz[baseline=-.5ex]\draw[message](0,0)--(.6,0);\ shared feedback: information influence\\[2pt]
\tikz[baseline=-.5ex]\draw[-{Latex[length=2mm]},draw=Gold,dashed,line width=.9pt](0,0)--(.6,0);\ shared evaluator: information influence\\[2pt]
\tikz[baseline=-.5ex]\draw[relation](0,0)--(.6,0);\ composition: interaction of selected patches\par
\vspace{.15cm}
A and B have different parents but share two sources of information.\par}
\end{columns}
\par\vspace{.05cm}\noindent\colorbox{Sand}{\parbox{\dimexpr\linewidth-2\fboxsep\relax}{\small{\color{Gold}\bfseries\fontsize{7.6}{9}\selectfont ANALOGY\,$\cdot$\,PHYSICS}\quad Our work on linker structure, branching and motor-driven avalanches motivated attention to connectivity. Stockmayer's gelation theory (1943) holds under its own assumptions. It does not predict an agent threshold.}}\par
\sources{Eliaz et al., PRE 2020. Liman et al., PNAS 2020. Li et al., JPCB 2021. \href{https://doi.org/10.1063/1.1723803}{Stockmayer, J.\ Chem.\ Phys.\ 11:45 (1943)}, analogy only.}
\note{[0:45] Ancestry alone forms a tree. Evidence flow forms a graph. Here A and B descend from different parents. Both children still read one feedback archive, and one evaluator scores both. The record therefore needs three edge meanings: state ancestry, information influence, and interaction between selected patches. Our biophysics work on branched networks drew our attention to connectivity. Stockmayer's gelation theory is only an analogy. Shared conditions can also help, as the next slide shows for paired comparisons.}
\end{frame}
